% sec_4_4_dynamics.tex — Динамика смазочного слоя

Разделы~4.1--4.3 рассматривали нелинейные эффекты применительно к~стационарному вращению вала при~фиксированном положении его~центра. В~реальных условиях эксплуатации вал совершает малые колебания относительно равновесного положения под~действием дисбаланса, переходных нагрузок и~гидродинамических возмущений. Реакция смазочного слоя на~эти~колебания описывается линеаризованными коэффициентами жёсткости и~демпфирования, которые определяют динамические свойства опорного узла и~устойчивость роторной системы в~целом.

\textit{Линеаризация по~малым возмущениям положения вала.}
Равновесное положение центра вала характеризуется координатами~$(x_0,\,y_0)$ или, эквивалентно, эксцентриситетом~$\varepsilon_0$ и~углом нагружения~$\psi_0$. При~малом возмущении центр вала смещается в~положение~$(x_0 + \xi,\; y_0 + \eta)$, где~$|\xi|,\,|\eta| \ll c$.

Безразмерная толщина зазора с~учётом возмущения записывается в~виде
\begin{equation}
  H(\theta,\,t) = H_0(\theta)
    - \bar{\xi}(t)\,\cos\theta
    - \bar{\eta}(t)\,\sin\theta,
  \label{eq:H_perturbed}
\end{equation}
\begin{conditions}
$H_0(\theta) = 1 + \varepsilon_0\cos\theta$ -- равновесный зазор радиального подшипника~\eqref{eq:journal_h};\\
$\bar{\xi} = \xi/c$,\; $\bar{\eta} = \eta/c$ -- безразмерные составляющие возмущения положения центра вала.
\end{conditions}

Давление разлагается по~малому параметру:
\begin{equation}
  P = P_0
    + \bar{\xi}\,P_{\xi}
    + \bar{\eta}\,P_{\eta}
    + \dot{\bar{\xi}}\,P_{\dot\xi}
    + \dot{\bar{\eta}}\,P_{\dot\eta}
    + \ldots
  \label{eq:P_expansion}
\end{equation}
Здесь~$P_0$~-- стационарное давление (главы~5--7); поля~$P_{\xi}$, $P_{\eta}$, $P_{\dot\xi}$, $P_{\dot\eta}$ находятся из~линеаризованных уравнений Рейнольдса с~правыми частями, получаемыми дифференцированием~\eqref{eq:reynolds_nondim} по~соответствующим переменным.

В~постановке JFO~\eqref{eq:jfo_conservation} граница кавитационной области при~возмущении может смещаться. В~первом приближении используется допущение о~фиксированной кавитационной области (\textit{frozen cavitation}): граница~$\partial\Omega_{\mathrm{cav}}$ определяется по~равновесному полю~$P_0$. Альтернативные подходы и~их~сравнение рассматриваются в~главе~8.

\textit{Матрицы жёсткости и~демпфирования смазочного слоя.}
Линеаризация гидродинамической силы приводит к~выражению
\begin{equation}
  \begin{pmatrix} \Delta F_x \\ \Delta F_y \end{pmatrix}
  = -\mathbf{K}
  \begin{pmatrix} \bar\xi \\ \bar\eta \end{pmatrix}
  - \mathbf{C}
  \begin{pmatrix} \dot{\bar\xi} \\ \dot{\bar\eta} \end{pmatrix},
  \label{eq:linearized_force}
\end{equation}
где~компоненты матриц жёсткости~$\mathbf{K}$ и~демпфирования~$\mathbf{C}$ определяются как
\begin{equation}
  K_{ij} = -\frac{\partial F_i}{\partial \bar{q}_j}\bigg|_0, \qquad
  C_{ij} = -\frac{\partial F_i}{\partial \dot{\bar{q}}_j}\bigg|_0,
  \quad i,j \in \{x,\,y\},
  \label{eq:KC_def}
\end{equation}
а~численное вычисление выполняется интегрированием возмущённых давлений по~области смазочного слоя:
\begin{equation}
  K_{xx} = -\iint_\Omega P_{\xi}\,\cos\theta\;d\theta\,dZ, \qquad
  K_{xy} = -\iint_\Omega P_{\eta}\,\cos\theta\;d\theta\,dZ,
  \label{eq:K_integrals}
\end{equation}
аналогично для~$K_{yx}$, $K_{yy}$ и~для~$C_{xx}$, $C_{xy}$, $C_{yx}$, $C_{yy}$ через~$P_{\dot\xi}$, $P_{\dot\eta}$.

Детерминированные микроструктуры изменяют стационарное поле~$P_0$ и~возмущённые поля~$P_{\xi}$, $P_{\eta}$, $P_{\dot\xi}$, $P_{\dot\eta}$, а~следовательно, компоненты~$K_{ij}$ и~$C_{ij}$. Это~является основным каналом влияния текстуры поверхности на~динамические характеристики опорного узла. Матрицы~$\mathbf{K}$ и~$\mathbf{C}$ в~общем случае несимметричны: перекрёстные коэффициенты~$K_{xy} \neq K_{yx}$, $C_{xy} \neq C_{yx}$, что~обусловлено вращательным характером течения смазки и~является принципиальным отличием подшипников скольжения от~упругих опор.

\textit{Критерии устойчивости роторно-подшипниковой системы.}
Уравнение движения ротора массой~$M$ в~двухкоординатной ($2$DOF) модели записывается в~виде
\begin{equation}
  M\,\ddot{\mathbf{q}} + \mathbf{C}\,\dot{\mathbf{q}} + \mathbf{K}\,\mathbf{q} = 0,
  \qquad \mathbf{q} = (\bar\xi,\,\bar\eta)^T.
  \label{eq:rotor_eom}
\end{equation}
Переход к~переменным состояния~$\mathbf{u} = (\mathbf{q},\,\dot{\mathbf{q}})^T$ приводит систему к~стандартной форме
\begin{equation}
  \dot{\mathbf{u}} = \mathbf{A}\,\mathbf{u}, \qquad
  \mathbf{A} =
  \begin{pmatrix}
    \mathbf{0} & \mathbf{I} \\
    -M^{-1}\mathbf{K} & -M^{-1}\mathbf{C}
  \end{pmatrix}.
  \label{eq:state_matrix}
\end{equation}
Система асимптотически устойчива тогда и~только тогда, когда все~собственные значения~$\lambda_k$ матрицы~$\mathbf{A}$ имеют отрицательную вещественную часть: $\mathrm{Re}(\lambda_k) < 0$ для~всех~$k$. Если~хотя~бы~для~одного~$k$ выполняется~$\mathrm{Re}(\lambda_k) > 0$, система неустойчива; граничный случай~$\max_k \mathrm{Re}(\lambda_k) = 0$ соответствует порогу устойчивости (критической скорости).

Формулы для~критических скоростей и~интегральных критериев устойчивости зависят от~конкретной модели ротора (масса, жёсткость вала, гироскопика) и~в~данном разделе не~приводятся; количественный пример применительно к~подшипнику центробежного насоса рассматривается в~главе~8.

\begin{figure}[!ht]
\centering
\fbox{\parbox{0.85\textwidth}{\centering
  \vspace{3.5cm}
  ПЛЕЙСХОЛДЕР ДЛЯ РИСУНКА\\
  Схема линеаризации: равновесное положение центра вала $(x_0, y_0)$,\\
  возмущение $(\xi, \eta)$, компоненты реакции смазочного слоя\\
  $\Delta F_x$, $\Delta F_y$ через матрицы $\mathbf{K}$ и $\mathbf{C}$
  \vspace{3.5cm}
}}
\caption{Линеаризованная модель сил смазочного слоя при малых колебаниях вала}
\label{fig:linearized_forces}
\end{figure}
